% !TeX root = ../nc-quadric-avatar.tex
% 2026-11-22 / cubic reachability and derived obstruction
\section{二重壁でのcubic型への到達不能性}
\label{sec:qa-as-cubic-obstruction}

二つの乗法がともに階数三となることだけでは、
他のmutation列や別のexceptional foundationによる
cubic型への到達を排除できない。
本節ではこの強い否定方向を
\cite[第19--21節]{AvatarVSeven2026}から
必要な数値・層論・平滑化の証明を保って統合する。
分類対象は引き続き全整数頂点を保つAS代数、
比較対象は構造対象を指定しない商Abel圏・導来圏である。

\subsection{有限foundationでのcubic候補の条件}
\label{subsec:qa-as-cubic-candidate}

$A\in\mathrm{AS}_{\mathrm{sq}}$の有限foundationを
$B=\End(E_0,E_1,E_2,E_3)$、
$\mathcal D=\operatorname{per}B\simeq D^b(\qgr A)$とする。
そのEuler行列は
\[
G_{\mathrm{sq}}=
\begin{pmatrix}
1&2&2&4\\0&1&0&2\\0&0&1&2\\0&0&0&1
\end{pmatrix}.
\]
\begin{definition}[cubic候補foundation]
\label{def:qa-as-cubic-candidate}
$\mathcal D$の四対象$(Q_0,Q_1,Q_2,Q_3)$が
cubic候補foundationであるとは、
strong exceptionalであり、
その$K_0(\mathcal D)$の類が整数基底をなし、
Euler行列が
\[
G_c=
\begin{pmatrix}
1&2&4&6\\0&1&2&4\\0&0&1&2\\0&0&0&1
\end{pmatrix}
\]
で、二つの隣接合成
\begin{equation}
\Hom(Q_{i+1},Q_{i+2})\otimes\Hom(Q_i,Q_{i+1})
\xrightarrow{\ \sim\ }\Hom(Q_i,Q_{i+2})
\quad(i=0,1)
\label{eq:qa-as-cubic-two-compositions}
\end{equation}
が同型であることをいう。
fullnessは別途仮定しない。
任意のcubic AS代数の標準四対象はこの条件を満たす。
\end{definition}

\begin{lemma}[cubic候補のrankと数値線束型]
\label{lem:qa-as-cubic-candidate-numerical}
cubic候補$(Q_0,\ldots,Q_3)$が存在すれば
共通の$[1]$ shiftの後、全$Q_i$の数値rankは一、
かつ適当な$a,b\in\mathbb Z$に対して
\[
c_1(Q_i)=(a,b),\ (a+1,b),\
(a+1,b+1),\ (a+2,b+1)
\]
または二つのrulingを交換した順序を取る。
数値Chern座標を$(r,a,b,s)$と書けば
各$Q_i$は$s=ab$を満たす。
\end{lemma}
\begin{proof}
squareの標準基底に対応する数値座標を
\[
[E_0]=(1,0,0,0),\quad [E_1]=(1,1,0,0),\quad
[E_2]=(1,0,1,0),\quad [E_3]=(1,1,1,1)
\]
とする。$x^2=y^2=0,\ xy=1$の形式的Chern記号で
Serre作用は
\[
S_*(r,a,b,s)=(r,a-2r,b-2r,s-2a-2b+4r)
\]
であり、
\[
(S_*-1)^2v=8r(v)\delta,\qquad
\delta=(0,0,0,1).
\]
cubicのEuler行列$G_c$から計算した
同じ作用の二乗は、四つの候補基底ベクトルを
同一のprimitiveな交代結合
$[Q_0]-[Q_1]-[Q_2]+[Q_3]$へ送る。
従って各rankは共通整数$r$であり、
この結合が$K_0$のprimitive元なので$r=\pm1$。
共通shiftで$r=1$とする。
自己Euler値は$1+2s-2ab=1$なので$s=ab$。
二つのrank一数値類のEuler値は
$(1+a_j-a_i)(1+b_j-b_i)$。
隣接方向のEuler値は二、逆向きは零だから、
差は$(1,0)$または$(0,1)$。
二段隔てたEuler値四により差は$(1,1)$となるので、
隣接差は交互となり主張を得る。
\end{proof}

\begin{lemma}[曲線への制限とdeterminant]
\label{lem:qa-as-cubic-restriction-determinant}
$(D,F,G,P)$をsquare avatarとし、
$\rho:\operatorname{per}B\to\Perf(D)$を
四線束$(P,F,G,FG)$へのrealizationとする。
任意の$X,Y\in\mathcal D$に自然な三角形
\begin{equation}
\RHom(Y,X)^\vee[-2]\longrightarrow
\RHom(X,Y)\longrightarrow
\RHom_D(\rho X,\rho Y)
\label{eq:qa-as-cubic-spherical-restriction}
\end{equation}
がある。
数値類$[X]=(r,a,b,s)$に対し
\begin{equation}
\det(\rho X)=F^aG^bP^{\,r-a-b+s}.
\label{eq:qa-as-cubic-determinant}
\end{equation}
$D$が滑らかで$X$がrank一のexceptional対象なら
$\rho X$は次数零のshiftを除き線束である。
\end{lemma}
\begin{proof}
前節のdg双加群三角形
\eqref{eq:qa-as-dg-triangle}とSerre双対性を
$\RHom(X,-)$へ適用すると
\eqref{eq:qa-as-cubic-spherical-restriction}を得る。
$X=Y$がexceptionalなら、
右辺の自己$\RHom$は$k\oplus k[-1]$となる。
$D$が滑らかなら$\coh D$は遺伝的で、
有界複体はコホモロジー層のshiftの直和になる。
Endが$k$でrankが一だから、
制限は一つのrank-one indecomposable vector bundleの
shift、すなわち線束のshiftである。
（torsionとvector bundleの拡大は分裂する。）
rank一の符号を共通shiftで正に揃える。

数値基底の座標を解くと
\[
[X]=(r-a-b+s)[E_0]+(a-s)[E_1]
 +(b-s)[E_2]+s[E_3].
\]
perfect複体のdeterminantの加法性と
$\rho(E_0,E_1,E_2,E_3)=(P,F,G,FG)$から
\eqref{eq:qa-as-cubic-determinant}を得る。
\end{proof}

\begin{proposition}[滑らかな二重壁では任意のcubic候補が不存在]
\label{prop:qa-as-cubic-smooth-no-candidate}
$D$を滑らかな楕円曲線とし、
$P^{\otimes2}\simeq\OO_D$、
$P\not\simeq\OO_D$、$G\simeq FP$とする。
$\mathcal D=\operatorname{per}B(D,F,G,P)$には
cubic候補foundationは存在しない。
\end{proposition}
\begin{proof}
候補を仮定する。
補題~\ref{lem:qa-as-cubic-candidate-numerical}により
全rankを一とし、補題
~\ref{lem:qa-as-cubic-restriction-determinant}により
$\rho Q_i=N_i[2m_i]$と書く。
隣接する$N_i$の次数差は二で、
$\Ext$の三角形と正次数線束の$H^1=0$から
前向きHomが次数零に集中するため、
各$m_i$は共通である。
共通shiftを除き、$N_i$を四線束とする。

\eqref{eq:qa-as-cubic-determinant}で
$G=FP$、$s=ab$、$P^2=\OO_D$を用いると
$\det(\rho Q)=F^{a+b}P^{1-a+ab}$となる。
数値差が$(1,0),(0,1),(1,0)$の場合、
連続する三つの商線束は
\[
(N_1N_0^{-1},N_2N_1^{-1},N_3N_2^{-1})
=(FP^{b-1},FP^{a+1},FP^b).
\]
逆順のrulingなら
$(FP^a,FP^b,FP^{a+1})$となる。
$P$が位数二なので商は$F,FP$の二種類だけで、
最初と最後が異なる。
従って三つのうち隣接する二つは同型である。
その共通次数二線束を$L$とすると、
合成写像
\[
 H^0(D,L)\otimes H^0(D,L)\longrightarrow H^0(D,L^2)
\]
のkernelは
$\bigwedge^2H^0(D,L)$という一次元空間である。
他方\eqref{eq:qa-as-cubic-spherical-restriction}によって
候補のforward Homと合成は
この線束上のHomと合成に一致する。
これは\eqref{eq:qa-as-cubic-two-compositions}の
同型性に矛盾する。
特定のmutation列に依存せず、
全てのcubic候補を排除した。
\end{proof}
